% Part: sets-functions-relations
% Chapter: functions
% Section: functions-relations

\documentclass[../../../include/open-logic-section]{subfiles}


\begin{document}

\olfileid{sfr}{fun}{rel}

\olsection{संबंध म्हणून फलने}

\begin{explain}
~$A$ मधील !!{element}s ची~$B$ मधील !!{element}s शी संगती लावणारे
फलन $A$ आणि~$B$ यांच्यामधील एक संबंध ठरवते: $f(x) = y$ असेल
तेव्हा आणि केवळ तेव्हाच $x$ आणि $y$ यांच्यात हा संबंध असतो.
खरे तर, उदाहरणार्थ गणिताच्या पायाभूत घटकांची संख्या कमी करायची असेल,
तर आपण फलन~$f$ आणि हा संबंध, म्हणजे जोड्यांचा एक संच, \emph{एकरूप}
मानू शकतो. मग असा प्रश्न उभा राहतो: कोणते संबंध अशा रीतीने फलने ठरवतात?
\end{explain}

\begin{defn}[फलनाचा आलेख] $f\colon A \to B$ हे फलन असू द्या.
~$f$ चा \emph{आलेख} म्हणजे पुढील व्याख्येने मिळणारा संबंध
$R_f \subseteq A \times B$:
\[
R_f = \Setabs{\tuple{x,y}}{f(x) = y}.
\]
\end{defn}

\begin{explain}
विस्तारात्मकतेमुळे फलनाचा आलेख एकमेव ठरतो. शिवाय, संचांसाठीच्या
विस्तारात्मकतेमुळे फलनांसाठीचे अप्रत्यक्ष विस्तारात्मकता तत्त्व लगेच
समर्थित होते: $f$ आणि~$g$ यांचा प्रांत व सहप्रांत एकच असेल आणि
सर्व मूल्यांवर ती जुळत असतील, तर ती एकच फलन असतात.

तसेच, एखादा संबंध ``फलनरूप'' असेल, तर तो एका फलनाचा आलेख असतो.
\end{explain}

\begin{prop}\ollabel{prop:graph-function}
$R \subseteq A \times B$ हा संबंध पुढील अटी पूर्ण करतो असे समजा:
\begin{enumerate}
\item $Rxy$ आणि $Rxz$ असल्यास $y = z$ असते; आणि
\item प्रत्येक $x \in A$ साठी $\tuple{x,
y} \in R$ असणारा काही $y \in B$ असतो.
\end{enumerate}
मग $f(x) = y$ तेव्हा आणि केवळ तेव्हाच $Rxy$, अशा व्याख्येने मिळणाऱ्या
फलन $f\colon A \to B$ चा आलेख $R$ असतो.
\end{prop}

\begin{proof}
$Rxy$ असणारा एक $y$ आहे असे समजा. $Rxz$ असणारा दुसरा $z \neq y$
असता, तर~$R$ वरील अट मोडली गेली असती. म्हणून $Rxy$ असणारा एक $y$
असेल, तर हा $y$ एकमेव असतो. त्यामुळे $f$ सुस्पष्टपणे परिभाषित आहे.
उघडच, $R_f = R$.
\end{proof}

\begin{explain}
प्रत्येक फलन $f\colon A \to B$ ला एक आलेख असतो, म्हणजे $f(x) = y$
ने परिभाषित होणारा A आणि B यांच्यातील, $A \times B$ मध्ये
समाविष्ट असलेला एक संबंध असतो. उलट,
\olref{prop:graph-function} मधील गुणधर्म असलेला प्रत्येक संबंध~$R
\subseteq A \times B$ हा एका फलन~$f \colon A \to B$ चा आलेख असतो.
फलने आणि त्यांचे आलेख यांचा हा निकटचा संबंध असल्यामुळे फलन म्हणजे
त्याचा आलेखच असे आपण मानू शकतो. दुसऱ्या शब्दांत, फलने विशिष्ट
संबंधांशी, म्हणजे विशिष्ट क्रमित घटकसमूहांच्या संचांशी, एकरूप मानता
येतात. \oliflabeldef{sfr:rel:ref:sec}{पण या ``एकरूप मानण्याचा'' आशय
\olref[sfr][rel][ref]{sec} मधल्यासारखाच आहे, हे लक्षात घ्या: तो
फलनांच्या सत्तामीमांसेविषयीचा दावा नाही, तर फलनांना विशिष्ट संच
\emph{मानणे} सोयीचे आहे हे निरीक्षण आहे. ही सोय होण्याचे एक कारण असे की, }{}
संबंधांवर जशा क्रिया केल्या तशाच क्रिया फलनांवर करण्याचा आता आपण
विचार करू शकतो (\olref[sfr][rel][ops]{sec} पाहा). विशेषतः:
\end{explain}
%READERNOTE{OLFUN-004 — इंग्रजी स्रोत येथे graph ला relation on A times B म्हणतो; पण त्याच ग्रंथाच्या व्याख्येनुसार त्याचा शब्दशः अर्थ (A times B) च्या वर्गाचा उपसंच होईल. मराठी मजकुरात अचूक अर्थ A आणि B यांच्यातील, A times B मध्ये समाविष्ट संबंध असा दिला आहे.}

\begin{defn}\ollabel{defn:funimage}
$f \colon A \to B$ हे फलन असू द्या आणि $C\subseteq A$ असू द्या.

~$f$ चे~$C$ पर्यंतचे \emph{मर्यादन} म्हणजे सर्व $x \in C$ साठी
$(\funrestrictionto{f}{C})(x) = f(x)$ या व्याख्येने मिळणारे
फलन~$\funrestrictionto{f}{C}\colon C \to B$. दुसऱ्या शब्दांत,
$\funrestrictionto{f}{C} = \Setabs{\tuple{x, y} \in R_f}{x \in C}$.

~$f$ चे~$C$ वरील \emph{उपयोजन} म्हणजे $\funimage{f}{C} =
\Setabs{f(x)}{x \in C}$. यालाच~$f$ अंतर्गत~$C$ ची \emph{प्रतिमा}
असेही म्हणतो.
\end{defn}

\begin{explain}
या व्याख्यांवरून कोणत्याही फलन~$f$ साठी $\ran{f} =
\funimage{f}{\dom{f}}$ असते.
\oliflabeldef{sfr:rel:ops:sec}{\olref[sfr][rel][ops]{sec} मधील
व्याख्या आणि फलनांना संबंधांशी एकरूप मानण्याची आपली पद्धत लक्षात
घेतल्यास या संकल्पना अनुरूप आहेत. मात्र फलनाचे C पर्यंतचे मर्यादन
केवळ आदान निर्देशांक C मध्ये ठेवते; संबंधाचे C वरील मर्यादन
दोन्ही निर्देशांक C मध्ये ठेवते. त्यामुळे त्या अगदी एकच क्रिया नाहीत.
व्युत्क्रम आणि सापेक्ष गुणाकार या इतर दोन क्रियांना थोडे अधिक
स्पष्टीकरण लागते. ते आपण
\olref[inv]{sec} आणि \olref[cmp]{sec} मध्ये देऊ.}{}
\end{explain}
%READERNOTE{OLFUN-005 — फलनाच्या मर्यादनाचे स्पष्ट सूत्र योग्य आहे: आलेखात पहिला निर्देशांक C मध्ये असतो. आधीचे संबंधावरील मर्यादन दोन्ही निर्देशांक C मध्ये ठेवते. इंग्रजी स्रोताचा exactly असा दावा येथे अनुरूप पण भिन्न असा स्पष्ट केला आहे.}

\end{document}
